\clearpage
\VerdierPhysicalPageAnchor{0196}%
\begingroup\small\setlength{\parskip}{0pt}\setlength{\emergencystretch}{1.5em}
\setlength{\abovedisplayskip}{4pt plus 2pt minus 2pt}%
\setlength{\belowdisplayskip}{4pt plus 2pt minus 2pt}%

\VerdierSemanticAnchor{ch3d-2-3-1-construction}
この構成ができれば、
$L^\bullet=\varprojlim_n L_n^\bullet$、
$w=\varprojlim_n w_n$ とおけばよい。$n$ が十分小さいとき、
$L_n=Y$、$v_n=\operatorname{id}_Y$、
$w_n=\operatorname{id}_Y$ とおく。すべての整数 $n<q$ に対し、
ここで $q\in\mathbf Z$ として、上の性質 1)、2)、3) をもつ複体
$L_n^\bullet$ と射 $v_n$、$w_n$ が構成されたと仮定する。複体 $X_n$ は次の形をする：

\VerdierSemanticAnchor{ch3d-2-3-1-xn-shape}
\[
X_n=\cdots\longrightarrow Y^{-n-2}\longrightarrow Z^{-n-1}
\longrightarrow X^{-n}\longrightarrow X^{-n+1}\longrightarrow\cdots
\]
また、複体 $L_n^\bullet$ は次の形をする：

\VerdierSemanticAnchor{ch3d-2-3-1-ln-shape}
\[
L_n^\bullet=\cdots\longrightarrow Y^{-n-2}\longrightarrow Z'^{-n-1}
\longrightarrow L^{-n}\longrightarrow L^{-n+1}\longrightarrow\cdots,
\]
ここで $Y^i$（それぞれ $X^i$）は複体 $Y$（それぞれ $X$）の成分である。
したがって、次の図式を得る：

\VerdierSemanticAnchor{ch3d-2-3-1-degree-q-diagram}
\begin{center}
\begin{tikzcd}[column sep=5.2em,row sep=2.8em]
  & X^{-q} \arrow[d,"(\dot u_q)^{-q}"] \\
Z'^{-q} \arrow[r,"(w_{q-1})^{-q}"] & Z^{-q},
\end{tikzcd}
\end{center}
ここで $(\dot u_q)^{-q}$ はエピ射であり、$Z'^{-q}\in M^\wedge$ である。
集合 $M^\wedge$ は性質 $(E_2)$（2.2.4）をもつので、この図式を次の可換図式に補える：

\VerdierSemanticAnchor{ch3d-2-3-1-lifting-square}
\begin{center}
\begin{tikzcd}[column sep=5.2em,row sep=2.8em]
L^{-q} \arrow[r,"\beta"] \arrow[d,"\alpha"']
  & X^{-q} \arrow[d,"(\dot u_q)^{-q}"] \\
Z'^{-q} \arrow[r,"(w_{q-1})^{-q}"] & Z^{-q},
\end{tikzcd}
\end{center}
ここで射 $\alpha$ はエピ射であり、対象 $L^{-q}$ は $M$ に属する。
(2.1.1.2) と (2.1.1.3) の記法を用いると、次の複体の可換図式が得られる：

\VerdierSemanticAnchor{ch3d-2-3-1-complex-square}
\begin{center}
\begin{tikzcd}[column sep=4.8em,row sep=3.0em]
L_{q-1}^\bullet(\alpha,-q) \arrow[r,"\varphi"]
  \arrow[d,"{p(\alpha,-q)}"']
  & X_{q-1}((\dot u_q)^{-q},-q)\simeq X_q
    \arrow[d,"{p((\dot u_q)^{-q},-q)\simeq\dot u_q}"] \\
L_{q-1}^\bullet \arrow[r,"w_{q-1}"] & X_{q-1}.
\end{tikzcd}
\end{center}
\endgroup
